\chapter{Continuum Boundary Radiation}
\label{chap:continuum-boundary-radiation}

\section{After Finite Forcing}

The previous chapter ended by forcing the continuum reading from finite
conditions. The ledger did not assume a completed line of points. It carried
spline atoms, finite supports, compatibility tests, extensions, and shrinking
bounds. A continuum reading was permitted only as the quotient of refinement
histories that no dense finite decision could still tell apart. That result is
powerful, but it is not yet the end of the grammar.

The next danger is subtler than the old temptation to assume the continuum in
advance. Once finite conditions have forced a continuum reading, the reader may
try to forget that the reading came from finite conditions at all. It may speak
as if completion erased the ledger that compelled it. It may also try to use
the completed space to repair what the finite split had forbidden: an interior
mixed scalar between the geometric face and the gauge face. This chapter is the
discipline that blocks those shortcuts.

The chapter begins by making the finite Cohen ledger extensional. Conditions
that carry the same atoms, supports, bounds, and decisions must be identified,
even when their presentations differ. That quotient must not add new
structure. It has to preserve finite support and countability. The continuum
reading is therefore lifted from the finite ledger, but the ledger remains
recoverable.

The second turn reads the gauge commutator residue. The residue does not close
because a completed algebra has been declared. It closes only because the
successor count can carry the ordered transports and their comparison. The
grammar must be able to count the first move, the second move, and the
residual difference between doing them in opposite orders.

The third turn opens the Sobolev--Hilbert door. The finite residue maps into
completion through an earned norm, and the completion preserves the split
between geometric and gauge readings. Completion does not mean fusion. It
means that Cauchy refinement has become stable enough to carry the same
obligations in a completed language.

The fourth turn removes the interior mixed term. Once the split has lifted,
compatibility forbids a continuum scalar that would multiply the geometric and
gauge channels inside the same interior certificate. The grammar does not lose
their relation. It refuses to put their reconciliation where the finite split
has already said it cannot live.

The fifth turn is the title's destination. If the interior mixed term is
forbidden and the terminal residue still has to be read, the boundary trace
must carry it. Radiation is therefore not an added image after the mathematics
has finished. It is the boundary phenomenon forced by the grammar: the residue
that the interior may not mix, lifted through completion and traced at the
boundary.

\section{The Extensional Ledger}

Finite Cohen forcing ended with conditions ordered by extension. A condition
carried a finite family of atoms, each with bounded support and a decision
made within the current tolerance. Extension preserved the old commitments and
added sharper ones. Compatibility forbade a later condition from erasing what
an earlier one had already made readable. That was enough to force a continuum
reading from below, but it left a bookkeeping question. When are two
conditions the same condition?

They cannot be the same merely because they were written the same way. A
finite measurement grammar is not allowed to attach truth to presentation.
The order in which atoms were listed, the name assigned to a support, or the
path by which two compatible fragments were merged cannot become new
continuum structure. If two conditions carry the same obligations, then for
the purposes of the ledger they must be identified.

The chapter therefore quotients the finite Cohen ledger extensionally. Two
conditions are extensionally the same when every atom carried by one is
carried by the other, every overlap agrees within the same bound, and every
decision forced by one is forced by the other. The quotient does not say that
all presentations are harmless. It says that presentation counts only when it
changes the finite obligations the condition carries.

In symbols, the section needs no more than the following reminder:
\[
  p \sim q
  \quad\hbox{when}\quad
  p\hbox{ and }q\hbox{ carry the same finite decisions.}
\]
The symbol is not the authority. The authority is the grammar that decides
what can be recovered by refinement. If a distinction between \(p\) and \(q\)
cannot be recovered by any finite atom, support, bound, or decision, then the
distinction is not a measurement distinction. It is quotiented away.

This quotient preserves the finite horizon. An extensional condition is still
finite because it is represented by finitely many carried atoms. It may have
many presentations, but it does not acquire infinitely many active
obligations. The quotient removes duplicate descriptions; it does not add a
hidden continuum behind the descriptions. Finite support remains the rule by
which the ledger stays honest.

Merge also becomes extensional. If two compatible conditions are merged, the
merged condition is read by the obligations carried in the union, not by the
route through which the union was assembled. Merging \(p\) with \(q\) and
then with \(r\) must identify with merging \(q\) with \(p\) and then with
\(r\), provided the same obligations survive. The ledger may remember the
refinement history when the history carries a real decision. It may not treat
mere order of inscription as a new decision.

Extension is read in the same way. A stronger condition extends a weaker one
when it carries every obligation of the weaker condition and adds a further
bounded decision. If a proposed strengthening abandons an old atom on shared
support, it is not an extension. If it merely writes the old atom under a
different name while preserving the same obligation, it is the same old
condition extensionally. The grammar identifies presentation and preserves
commitment.

This is where countability matters. The finite ledger can be counted because
its atoms can be indexed by finite data: support, bound, local value, and the
small list of neighboring obligations against which compatibility is checked.
A finite condition is then a finite list of finitely indexed atoms. Finite lists
of countable data remain countable. Extensional quotienting can only identify
some of those lists with one another. It does not create a larger collection.
The countable skeleton survives.

The Cantor--Godel--Cohen warning illuminates the point if used carefully.
Completions can differ by distinctions that no chosen finite ledger is able
to recover. The present grammar does not deny that such stories can be told.
It denies them the status of forced measurement structure. A completion is
admissible here only where its distinctions can be pressed back down into
finite refinements. Extensionality is the rule that keeps the upward
completion answerable to the downward record.

Precision supplies a second illustration. A spline completion may assign an
intermediate value between anchors. That value is not a new measurement, and
it is not an arbitrary guess. It is a consequence of the admissible completion
carried by the anchors. In the same way, an extensional condition may have a
completed reading, but the completed reading is answerable to the finite
obligations that forced it. If two presentations force the same obligations,
their completed readings are identified.

The extensional ledger is therefore the first lift of the chapter. Finite
conditions become extensional without ceasing to be finite. The quotient
identifies redundant presentations, preserves finite support, and carries
countability through the continuum door. Once the ledger has become
countable in this stricter sense, the next residue can be read. The gauge
commutator will ask not merely what a condition carries, but what order of
transport the count is able to close.

\section{Successor Residue}

The extensional ledger says which finite obligations count as the same. It
does not yet say what happens when two gauge transports are both admissible
but do not commute. Gauge order matters because a transport is not a decorative
arrow between completed points. It is a rule for carrying a distinction from
one finite reading to the next. If the order of two transports changes the
carried distinction, the difference is a residue the grammar must read.

The usual notation for the commutator is short:
\[
  [A,B] = AB - BA .
\]
The notation is useful only after the grammar has earned it. The expression
\(AB\) says that the \(A\)-transport is carried and then the \(B\)-transport
is carried. The expression \(BA\) says the reverse. The bracket records the
debt left when those two ordered readings fail to identify. The residue is
not an algebraic ornament. It is the finite ledger's witness that order has
measurement content.

Why should this close only through successor counting? Because the grammar
cannot read an ordered difference without counting the order. The first
successor carries the first admissible transport. The second successor
carries the next transport. A further comparison reads whether the two-step
path can be identified with the path taken in reverse. Without the successor
discipline, the reader would speak as if \(AB\) and \(BA\) already lived in a
completed arena where subtraction could be applied freely. The grammar has
not granted that arena.

Successor counting is austere. It does not bring in all numbers at once. It
permits a next step, then another next step, and requires each step to say
what it carries. The commutator residue needs exactly that discipline. One
step alone cannot expose noncommutation. Two steps can create the ordered
comparison. The return comparison asks whether the two ordered readings close
on the same extensional condition. If they do not, the residue is the
unclosed debt.

The count also prevents a false cancellation. If the ledger forgot which
transport occurred first, \(AB\) and \(BA\) could be collapsed by presentation
instead of by measurement. Extensionality would then be misused. The
extensional quotient identifies conditions with the same carried obligations;
it does not erase order when order changes what is carried. The successor
count keeps the order recoverable, and therefore keeps the residue honest.

Closure means that the residue has been counted, bounded, and assigned to the
ledger. It does not mean the residue disappears by decree. A gauge
commutator closes when the grammar has enough successor structure to read
the two ordered paths, compare them extensionally, and carry any remaining
difference as a bounded residue. The close-out is finite even when it points
toward completion. It says: this is the residual debt produced by these
ordered transports, and this debt is now in the count.

A loop illustration can help. In a Mach--Zehnder reading, two admissible
branches may remain separately consistent until they meet at a shared
boundary. If their carried distinctions agree, the boundary fold identifies
them. If they return opposed, the fold cannot form a single ledger without
cancelling the inadmissible ordering. The example is optical in its familiar
setting, but the grammatical lesson is simpler: ordered branches can carry a
residue that is invisible to each branch alone and readable only at the
comparison.

The Aharonov--Bohm illustration says the same thing in gauge language. Local
segments may look silent with respect to the field reading, while the closed
transport around a loop returns a holonomy. The chapter does not need the
physical apparatus to establish the point. The point is grammatical: a
connection can carry memory around a comparison, and the commutator residue
is the finite mark that order has not quotiented away.

This is also why the count reaching three matters. The first count permits a
transport to be carried. The second count permits the ordered pair of
transports. The third count permits the comparison to return as a closed
ledger entry rather than as an untracked alternation. The grammar is not
celebrating the numeral. It is recording the minimum successor discipline
needed for a loop residue to become a stable finite obligation.

Once the commutator residue has closed in this way, it can be handed to
completion. But completion must not wash out its finite origin. The residue
was not born in a smooth space. It was born from extensional conditions and
successor-ordered transport. The next section therefore asks how this finite
residue maps into the Sobolev--Hilbert completion while preserving the split
that made the gauge reading lawful.

\section{The Sobolev--Hilbert Trace}

The residue cannot enter completion as a naked symbol. Completion is earned
only when the grammar has a norm that reads activity, detects zero, and
controls Cauchy refinement. Earlier chapters grounded that norm in the
positive invariant: energy squared became a readable activity, and the
square-root reading opened the Hilbert door. Here the same door must receive
the extensional Cohen ledger and the successor-closed commutator residue.

The finite norm has a simple obligation. It must read zero exactly when the
residue has no activity left to carry. In schematic form,
\[
  \|r\| = 0 \quad\hbox{only when}\quad r \hbox{ is the zero residue.}
\]
The formula is not a new foundation. It records the old boundary-anchored
positivity in the language needed for completion. A norm that allowed
nonzero residue to pass as zero would reopen the null drift that the
concrete operator had already killed.

The pre-Hilbert skeleton is built from the countable extensional ledger.
Its elements are finite-supported readings of the obligations already
carried: atoms, bounds, decisions, and residues. Addition and scaling are
allowed only insofar as they preserve the finite compatibility discipline.
The skeleton is not yet the completed space. It is the countable stage on
which Cauchy refinement can be read.

Completion begins when sequences in this skeleton become Cauchy under the
norm. A sequence is Cauchy when late refinements no longer differ by more
than the chosen bound. This is the same finite approach that has guided the
book from the beginning: the grammar does not claim to have traversed every
point, but it can say that every later admissible refinement is being squeezed
within the bound already named. The Hilbert completion is the quotient of
such Cauchy refinement histories by unreadable late-stage difference.

The finite residue maps into this completion by remembering its origin. A
completed residue is not a residue without ancestry. It is the completed
reading of a finite residual history whose stages remain recoverable from
the ledger. The map into completion therefore carries the residue; it does
not dissolve it. If a completed object cannot be traced back to finite
residual obligations, it is not admitted by this grammar.

This is the first sense of trace in the section. A trace is the readable
memory of the finite source inside the completed reading. The trace does not
make the completion smaller. It makes it answerable. When a completed
residue is challenged, the grammar can ask which finite supports, bounds,
and successor comparisons forced it. If none can be given, the alleged
completion is optional structure, not a forced reading.

The second sense of trace prepares the boundary. Sobolev language is useful
because it knows how an interior activity can have a boundary reading. But
the chapter must use that language only after the finite discipline is in
place. The boundary trace is not a gift from smooth analysis. It will be the
completed reading of terminal finite residue. The present section builds the
door through which that terminal residue can pass.

Most important, the geometry/gauge split survives the passage. The geometric
certificate family carries boundary separation, mesh refinement, and
curvature-like residue. The gauge certificate family carries transported
orientation, charge, and commutator residue. Completion may complete both
families, but it may not identify their channels. The finite split lifts as
orthogonality in the Sobolev--Hilbert reading.

In symbols, the discipline has the shape
\[
  G \perp Q,
\]
where \(G\) names the completed geometric channel and \(Q\) names the
completed gauge channel. The symbol \(\perp\) is again only a record. It
says that the completed inner product respects the finite carrier split. It
does not say that geometry and gauge are unrelated. It says that their
relation is not an interior scalar product between two freely summable
components.

Precision again illustrates the distinction. A spline completion may supply
intermediate values that were not observed, but those values remain
consequences of the anchor grammar. They do not become independent
measurements. Likewise, the Sobolev--Hilbert completion supplies completed
residue readings, but those readings remain consequences of finite
conditions, successor residue, and the norm that bounds refinement.

The section therefore lifts finite residue without betraying finitude. The
countable extensional ledger enters a pre-Hilbert skeleton. The norm detects
zero activity and reads Cauchy refinement. The finite residue maps into
completion with its trace intact. The geometric and gauge certificate
families complete as split channels. Once that split has lifted, the next
attempted shortcut is visible: an interior mixed scalar in the continuum.
The grammar must remove it.

\section{No Interior Mixed Term}

Completion is often where forbidden mixtures try to return. The finite
chapter had already split the geometric and gauge carriers. A finite
certificate could read the geometric face or the gauge face, but it could
not form an interior direct sum that smuggled in a scalar mixed term. The
Sobolev--Hilbert lift preserves that split. Therefore the completed grammar
must also forbid the interior product that would undo it.

The issue is not whether geometry and gauge are connected. The previous
chapter established the connection sharply: boundary geometry generated the
gauge. The issue is where the connection may be read. Since the gauge arose
because the interior could not settle the charged residue on its own, the
completed interior may not later claim a scalar certificate that multiplies
geometric and gauge residues as if both were freely present there.

Compatibility supplies the rule. A continuum reading is compatible only if
it preserves the obligations that finite refinement forced. It must preserve
extensional countability, successor residue, the normed completion, and the
geometry/gauge split. A smooth-looking expression that violates one of those
obligations is not a deeper reading. It is an inadmissible presentation.

Thus the interior mixed term cancels:
\[
  \langle g,q\rangle_{\mathrm{int}} = 0 .
\]
This is not a claim that the symbols \(g\) and \(q\) have no relation. It is
the statement that their interior scalar cross-reading is forbidden by
compatibility. The geometric channel \(g\) and the gauge channel \(q\) may
meet at the boundary process that generated the gauge, but the completed
interior cannot multiply them into a hidden scalar debt.

The cancellation has a grammatical source. The split lifted as
orthogonality. An interior mixed scalar would identify what orthogonality was
forced to keep apart. It would let the geometric channel pay gauge debt or
the gauge channel pay geometric debt inside the interior. That would erase
the boundary origin of the gauge and turn the forced split into a notation.
Compatibility forbids that erasure.

The interior is therefore poorer than the completed notation might suggest.
It can carry completed geometric readings. It can carry completed gauge
readings. It can carry the normed residue histories that led to both. But it
cannot carry an independent scalar product between the channels. The missing
term is not a gap. It is the discipline that keeps the completed reading
faithful to the finite one.

A Casimir-style boundary imbalance gives a useful illustration. When a
boundary restricts admissible interior refinements, the repair cannot be
made by inventing modes inside the restricted region. The repair is read at
the boundary where the imbalance is exposed. The example does not prove the
chapter's claim, but it shows the same grammatical shape: an interior
deficit is not healed by adding unearned interior structure.

The holographic illustration says the same thing at a larger scale. If every
admissible interior refinement must be reconciled through the boundary, then
volume is not a storehouse of independent record. Interior presentations are
smooth shadows of boundary bookkeeping. The present chapter is narrower: it
does not need a general scaling law. It needs only the rule that the
geometry/gauge mixed scalar has no independent interior certificate.

This also clarifies the role of compatibility in the continuum. Compatibility
is not passive agreement between completed objects. It is the active refusal
to let completion forget how the objects were earned. If a completed reading
would require an interior mixed term, then it fails compatibility with the
finite split. The grammar cancels the term before it can become a false
reconciliation.

The residue, however, has not vanished from the story. Removing the interior
mixed scalar does not remove the terminal debt that the geometric and gauge
readings jointly expose. It only says where that debt cannot be paid. The
residue must still be read, and the only remaining licensed place is the
boundary trace. The next section therefore forces radiation to the boundary.

\section{Radiation At The Boundary}

The boundary trace is the reader of terminal residue. A finite residual
history may refine through many interior stages, but when the grammar asks
where the unabsorbed debt is finally exposed, it reads the boundary. The
trace does not invent a new residue. It agrees with the terminal residue that
the interior completion has carried but is forbidden to mix as an interior
scalar.

The relation can be written schematically:
\[
  \operatorname{Tr}_{\partial}(r) = r_{\mathrm{term}} .
\]
The left side names the boundary trace. The right side names the terminal
residue of the finite-to-completed history. The equality says that the
boundary reading is not optional decoration. It is the authorized reading of
the residue left after compatible completion has removed the interior mixed
term.

The Sobolev--Hilbert completion matters here because it prevents the boundary
trace from being a merely finite endpoint. Residual histories are Cauchy in
the norm. Their late refinements become indistinguishable inside the chosen
bound, and the completion identifies those late-stage readings. The trace
then reads the completed terminal residue. It is finite in ancestry and
completed in form.

Radiation is the name for that boundary phenomenon. It is not introduced as
an independent substance moving through an already completed space. It is
forced when four obligations meet: the extensional ledger preserves finite
countability; successor counting closes the commutator residue; the
Sobolev--Hilbert door lifts the residue while preserving the split; and
compatibility forbids the interior scalar mixed term. What remains must be
carried by the boundary trace.

This is why the chapter's title joins continuum, boundary, and radiation.
The continuum is not assumed. It is forced from finite conditions and lifted
through completion. The boundary is not a decorative edge. It is the place
where terminal residue remains readable after the interior has been
forbidden to mix the channels. Radiation is not an afterthought. It is the
readable boundary form of the unabsorbed residue.

A Hawking-style illustration can help if it stays modest. A paired reading
near a horizon is not accepted as two independent interior facts. Once one
side is fixed by an outside reading, the other side is forced into a
complementary role across the boundary. The familiar physical story speaks
of outward radiation and inward partner. The grammar reads a more general
shape: boundary asymmetry turns an uncommitted pair into a traced residue.

The Casimir illustration gives another angle. A restricted interior may
carry fewer admissible refinements than its exterior. The repair is then not
an invented interior mode but a boundary correction that restores global
distinguishability. The chapter's radiation reading has the same form at a
more structural level. The interior mixed term is forbidden, so the repair
is read where compatibility still permits it: at the boundary.

The Mach--Zehnder illustration supplies a third lantern. Two branches remain
admissible until recombination forces a boundary comparison. Agreement
folds them into one ledger; opposition cancels the inadmissible ordering.
Again, the example is not the proof. It shows why a residue can be carried
silently through interior segments and become readable only at the boundary
where the grammar must reconcile the histories.

The boundary trace also preserves the geometry/gauge split. It does not
collapse the two channels into an interior sum. It reads their terminal
relation from the boundary, where the gauge was generated and where the
residue can be reconciled without violating compatibility. The split is
therefore not a problem radiation solves. It is the reason radiation is
forced to the boundary.

This completes the lift begun by the extensional ledger. Finite conditions
became extensional and countable. The commutator residue closed through
successor counting. The residue mapped into completion through an earned
norm. The split survived as Sobolev--Hilbert orthogonality. The interior
mixed scalar was removed. The terminal residue then had one lawful reader:
the boundary trace. Boundary radiation is the name of that forced reading.

\section{What The Boundary Carries Forward}

The chapter has not added a new primitive to the grammar. It has protected
the forced continuum reading from two temptations: unearned completion and
interior mixing. The extensional ledger identified duplicate presentations
without losing finite support. Countability survived because the quotient
only removed redundant finite indices. The continuum therefore remained
answerable to finite conditions.

The chapter then made order readable. The gauge commutator residue closed
only when successor counting could carry the ordered transports and their
comparison. That residue entered completion through the normed
Sobolev--Hilbert door. The finite origin was traced inside the completed
reading, and the geometry/gauge split lifted rather than dissolved.

Compatibility then did its most severe work. It forbade the completed
interior from forming a scalar mixed term between the geometric and gauge
channels. The term vanished not because the channels had no relation, but
because their relation was generated at the boundary and could not be
reinserted as an interior product without undoing the split.

Finally, the boundary trace read the terminal residue. Radiation was forced
there because the interior had no licensed mixed certificate left. The
boundary restored a place where the completed residue could be read without
violating the finite grammar that produced it.

The next chapter receives this boundary as an outside reader. Once radiation
has been forced to the boundary, the interior invariant can be read from
outside as an obstruction rather than as an interior possession. The
continuum has been completed, but completion has not silenced the ledger.
The boundary still carries the residue the interior cannot spend.
